
\documentclass[11pt,a4paper]{article}
\usepackage[spanish]{babel}
\usepackage[utf8]{inputenc}
\usepackage[T1]{fontenc}
\usepackage{amsmath,amssymb}
\usepackage{siunitx}
\usepackage{geometry}
\geometry{margin=2.5cm}

\begin{document}

\title{Problema p.2 -- Flujo de nitrógeno entre dos tanques\\
Modelo adiabático subsónico desde reservorio}
\author{Resolución detallada}
\date{}
\maketitle

\section*{1. Enunciado (resumen)}

Dos tanques idénticos A y B (10 m$^3$ cada uno) contienen nitrógeno a la misma temperatura ambiente (\SI{15}{\celsius}).
Se conectan mediante una conducción de acero comercial Sch~40 de 1/2'' de diámetro nominal y longitud real \SI{15}{\meter}, que incluye:
\begin{itemize}
  \item 5 codos estándar de \ang{90},
  \item 2 válvulas esféricas,
  \item 1 entrada con resalte hacia el interior.
\end{itemize}

Inicialmente:
\[
P_A = \SI{8.0}{bar(g)} \Rightarrow P_{0} = \SI{9.0}{bar(abs)},\qquad
P_B = \SI{4.0}{bar(g)} \Rightarrow P_{2} = \SI{5.0}{bar(abs)}.
\]

Se abren las válvulas y se permite el flujo hasta que las presiones se igualan.
Se pide: \emph{máximo flujo másico} que puede establecerse entre los tanques.
(Respuesta de referencia: $\dot m_{\max} \approx \SI{0.088}{kg/s}$).

\section*{2. Datos, incógnita y propiedades}

Gas: nitrógeno, modelado como gas ideal:
\[
M = \SI{28}{g/mol}, \qquad
R = \frac{R_u}{M} \approx \frac{8.314}{0.028} \simeq \SI{2.976\times 10^{2}}{J/(kg\,K)},
\]
\[
\gamma = \frac{c_p}{c_v} \approx 1.4.
\]

Temperatura inicial (y ambiente):
\[
T_0 = T_A = T_B = \SI{15}{\celsius} \approx \SI{288}{K}.
\]

Presiones absolutas iniciales:
\[
P_0 = \SI{9}{bar} = \SI{9\times10^5}{Pa}, \qquad
P_2 = \SI{5}{bar} = \SI{5\times10^5}{Pa}.
\]

Relación de presiones:
\[
F = \frac{P_2}{P_0} = \frac{5}{9} \simeq 0.5556.
\]

Tubería:
\begin{itemize}
  \item Diámetro interno 1/2'' Sch~40:
  \[
  D \simeq \SI{0.0158}{m}.
  \]
  \item Longitud real $L = \SI{15}{m}$.
  \item Longitud equivalente total (15 m + accesorios, usando Crane para 5 codos, 2 válvulas esféricas y entrada con resalte):
  \[
  L'_e \simeq \SI{18.5}{m}.
  \]
  \item Tubería de acero comercial, régimen completamente turbulento: factor de Darcy tomado de Crane para 1/2'':
  \[
  f_T = 0.026.
  \]
\end{itemize}

Volumen específico en el reservorio A:
\[
v_0 = \frac{R T_0}{P_0}
      \approx \frac{2.976\times10^{2} \cdot 288}{9\times10^5}
      \simeq 9.50\times10^{-2}\ \text{m}^3/\text{kg}.
\]

Incógnita principal: máximo flujo másico $\dot m_{\max}$.

\section*{3. Hipótesis físicas y elección del modelo ideal}

\begin{enumerate}
\item \textbf{Flujo estacionario cuasi-permanente.}
El máximo flujo ocurre al comienzo (cuando $P_A = \SI{9}{bar}$, $P_B = \SI{5}{bar}$).
Aproximamos ese instante por un flujo estacionario entre dos reservorios a presiones fijas $P_0, P_2$.

\item \textbf{Gas ideal, propiedades constantes} $(R,\gamma)$.

\item \textbf{Conducción horizontal de sección constante}, sin trabajo de eje:
modelo de ducto 1D con fricción.

\item \textbf{Modelo térmico y comparación con casos reales.}
Los tanques están inicialmente a \SI{15}{\celsius}, igual temperatura que el ambiente.
La conducción no está altamente aislada, pero tampoco tiene un intercambio intenso de calor.
Del análisis comparativo de modelos (documento S4):

\begin{itemize}
  \item Para $T_1 \approx T_{\text{amb}}$ y calor moderado, el flujo real se ubica entre los modelos
  isotérmico y adiabático:
  \[
    \left(\frac{w}{A}\right)_{\text{iso}}
    < \left(\frac{w}{A}\right)_{\text{real}}
    < \left(\frac{w}{A}\right)_{\text{adiab}}.
  \]
  \item Como se pide \emph{máximo flujo másico}, usamos el modelo \textbf{adiabático} como
  \textbf{cota superior razonable}.
\end{itemize}

Trabajaremos, por tanto, con el modelo \emph{adiabático con fricción, subsónico, desde reservorio}.

\item \textbf{Régimen subsónico.}
Verificaremos a posteriori que el número de Mach es suficientemente bajo;
resultará $M_2 \approx 0.22$, consistente con la hipótesis.
\end{enumerate}

\section*{4. Variables adimensionales y definición de $X$}

Seguimos la reducción de las notas para ``flujo adiabático desde un reservorio, subsónico''.

Definimos las variables adimensionales:
\[
r = \frac{v_0}{v_{1'}}, \qquad
z = \frac{v_0}{v_2},
\]
donde $1'$ es la sección de entrada efectiva del ducto y $2$ la de salida hacia el segundo tanque.

La relación de presiones es
\[
F = \frac{P_2}{P_0} = \frac{5}{9}.
\]

Definimos el parámetro de flujo
\[
X = \left(\frac{\dot m}{A}\right)^2 \frac{v_0}{P_0}, \qquad
A = \frac{\pi D^2}{4}.
\]

Introducimos además
\[
C = \frac{2\gamma}{\gamma-1}.
\]
Para $\gamma = 1.4$:
\[
C = \frac{2\cdot 1.4}{0.4} = 7.
\]

\section*{5. Sistema reducido de ecuaciones adiabáticas}

\subsection*{5.1. Ecuaciones reducidas (caso desde reservorio)}

De las transparencias para ``entrada desde reservorio estanco -- flujo adiabático'' el sistema es:

\begin{enumerate}
\item Relación reserva--entrada--salida:
\[
E_1(r,z) = r^2 - r^{\gamma+1} - z^2 + F z = 0.
\]

\item Definición de $X$ (por eliminación):
\[
X = C\,(z^2 - F z).
\]

\item Ecuación de Fanno reducida entre $1'$ y $2$:
\[
\frac{r^2 - z^2}{X} - \frac{\gamma+1}{\gamma} \ln\!\left(\frac{r}{z}\right)
   = f\, \frac{L'_e}{D}.
\]
\end{enumerate}

Aquí $f$ es el factor de fricción de \emph{Fanning} en la deducción original.
Disponemos del Darcy $f_T$ de Crane.
Con $f_{\text{Darcy}} = 4 f_{\text{Fanning}}$, las fórmulas reducidas se han reescrito de modo que el término de fricción a la derecha puede expresarse directamente como
\[
\text{RHS} = f_T \frac{L'_e}{D}.
\]

Sustituyendo $X = C(z^2 - F z)$ en la ecuación de Fanno, el sistema se escribe explícitamente como
\begin{align*}
E_1(r,z) &=
r^2 - r^{\gamma+1} - z^2 + F z = 0, \\
E_2(r,z) &=
\frac{r^2 - z^2}{C\,(z^2 - F z)}
- \frac{\gamma+1}{\gamma} \ln\!\left(\frac{r}{z}\right)
- f_T\frac{L'_e}{D} = 0.
\end{align*}

Con los datos numéricos:
\[
F=\frac{5}{9},\quad
\gamma=1.4,\quad
C=7,\quad
f_T=0.026,\quad
L'_e=\SI{18.5}{m},\quad
D=\SI{0.0158}{m},
\]
\[
\text{RHS} = f_T \frac{L'_e}{D}
= 0.026\;\frac{18.5}{0.0158}
\simeq 30.44.
\]

\section*{6. Resolución numérica del sistema $(r,z)$}

Las condiciones físicas para la solución subsónica son:
\begin{itemize}
  \item $0<z<r<1$,
  \item $z^2 - F z > 0$ (para que $X>0$).
\end{itemize}

Se usa un método de Newton--Raphson bidimensional, partiendo de semillas compatibles con el criterio de las notas:
\[
r^{(0)}\approx 0.97, \qquad
z^{(0)}\approx F+0.01\simeq 0.565.
\]

El algoritmo itera hasta que $|E_1|,|E_2| < 10^{-10}$.
La solución convergente que satisface las restricciones físicas es
\[
\boxed{
r \simeq 0.9923,\qquad
z \simeq 0.5610.
}
\]

Con estos valores:
\[
X = C\,(z^2 - F z)
   = 7\,[z^2 - F z]
   \simeq 0.02132.
\]

\section*{7. Cálculo del flujo másico $\dot m$}

Por definición,
\[
X = \left(\frac{\dot m}{A}\right)^2 \frac{v_0}{P_0}
\quad\Rightarrow\quad
\frac{\dot m}{A} = \sqrt{\frac{X P_0}{v_0}}.
\]

\subsection*{7.1. Volumen específico en A}

\[
v_0 = \frac{R T_0}{P_0}
    \simeq \frac{2.976\times10^{2} \cdot 288}{9\times10^5}
    \simeq 9.50\times10^{-2}\ \text{m}^3/\text{kg}.
\]

\subsection*{7.2. Área interna del tubo}

\[
A = \frac{\pi D^2}{4}
  = \frac{\pi (0.0158)^2}{4}
  \simeq 1.96\times 10^{-4}\ \text{m}^2.
\]

\subsection*{7.3. Flujo másico por unidad de área}

\[
\left(\frac{\dot m}{A}\right)
= \sqrt{\frac{X P_0}{v_0}}
= \sqrt{\frac{0.02132 \cdot 9\times10^5}{9.50\times10^{-2}}}
\simeq 4.49\times10^{2}\ \frac{\text{kg}}{\text{m}^2\text{s}}.
\]

\subsection*{7.4. Flujo másico total}

\[
\dot m = A\;\frac{\dot m}{A}
\simeq (1.96\times10^{-4})
       (4.49\times10^{2})
\simeq 8.8\times10^{-2}\ \text{kg/s}.
\]

Es decir,
\[
\boxed{\dot m_{\max} \approx \SI{0.088}{kg/s}}
\]
en perfecto acuerdo con la respuesta de la guía.

\section*{8. Verificación física}

\subsection*{8.1. Chequeo de unidades}

En
\[
X = \left(\frac{\dot m}{A}\right)^2 \frac{v_0}{P_0}
\]
tenemos:
\begin{itemize}
  \item $[\dot m/A] = \text{kg\,m}^{-2}\text{s}^{-1}$,
  \item $[v_0] = \text{m}^3\text{kg}^{-1}$,
  \item $[P_0] = \text{N\,m}^{-2} = \text{kg\,m}^{-1}\text{s}^{-2}$.
\end{itemize}
Entonces:
\[
\left[\left(\frac{\dot m}{A}\right)^2 \frac{v_0}{P_0}\right]
= \frac{\text{kg}^2\text{m}^{-4}\text{s}^{-2}\cdot \text{m}^3\text{kg}^{-1}}
        {\text{kg\,m}^{-1}\text{s}^{-2}}
= 1\quad (\text{adimensional}).
\]
Correcto.

\subsection*{8.2. Velocidad y número de Mach en la salida}

Volumen específico en 2:
\[
v_2 = \frac{v_0}{z} \simeq \frac{9.50\times10^{-2}}{0.561}
      \simeq 1.69\times10^{-1}\ \text{m}^3/\text{kg}.
\]

Flujo másico por unidad de área:
\[
\frac{\dot m}{A} \simeq 4.49\times10^{2}\ \text{kg/(m}^2\text{s)}.
\]

Densidad en 2: $\rho_2 = 1/v_2 \simeq 5.9\ \text{kg/m}^3$.

Velocidad en 2:
\[
u_2 = \frac{\dot m/A}{\rho_2}
      \simeq \frac{4.49\times10^{2}}{5.9}
      \simeq \SI{76}{m/s}.
\]

Para el nitrógeno a unos $\sim \SI{285}{K}$, la velocidad del sonido es del orden de
\[
a_2 = \sqrt{\gamma R T_2} \sim \SI{340}{m/s},
\]
por lo que
\[
M_2 = \frac{u_2}{a_2} \approx \frac{76}{344}\simeq 0.22 \ll 1.
\]

El flujo es claramente \textbf{subsónico}, coherente con el modelo adoptado.

\section*{9. Comentario final sobre el modelo}

La instalación real no es estrictamente adiabática, pero:
\begin{itemize}
  \item $T_1 \approx T_{\text{amb}}$,
  \item la conducción no es extremadamente larga,
  \item el enunciado pide un \emph{máximo flujo másico}.
\end{itemize}

Según el análisis comparativo de modelos:
\begin{itemize}
  \item Para este tipo de situación, el flujo real debería ubicarse entre el modelo isotérmico y el adiabático,
  \item El modelo adiabático funciona como una \textbf{cota superior razonable} de $(w/A)_{\text{real}}$.
\end{itemize}

El valor obtenido
\[
\dot m_{\max} \approx \SI{0.088}{kg/s}
\]
puede interpretarse como una estimación del \emph{máximo flujo másico instantáneo} al inicio del proceso,
físicamente consistente y en acuerdo con la solución de la guía.

\end{document}
